% a06_BALtorus_ct6: Step 3(F), second part: the quadratic-form step Ba2 (2.108)-(2.109) (the half of the stage of
% Ba2 Prop. 2.6 that concerns no interaction term), and the passage from Ba2 (3.5) at k = 0 to (3.21) and (3.42).

\subsubsection*{The quadratic-form step \cite[(2.108)--(2.109)]{Ba2}}

\emph{Setting.} Step $k$ in scale-$k$ units, $t\defeq t_k$, $\eta\defeq L^{-k}$; $\Omega_A\defeq B^k(\Lambda^{(k-1)'}_2)$ (the propagator domain of \cite[(2.53), (2.56)]{Ba2}), $\Omega_B\defeq B^k(\Lambda^{(k-1)'}_2\cap\Lambda^{(k)c}_5)$ and $\Omega_C\defeq B^k(\Lambda^{(k)}_2)$ (the domains of \cite[(2.108)]{Ba2}), all unions of big blocks; $\tilde B\defeq A^{(k)}$ on $\Omega_A$. On $\Omega_C$, $\tilde B=B^{(k+1),\eta}$: $\theta_{k+1}=1$ on $B^k(\Lambda^{(k)}_2)$ \cite[pp.~567, 576]{Ba2}, $\Omega_C\subset B^k(\Lambda^{(k)}_0)$, and by \cite[p.~579]{Ba2} the configuration there is $(1-\theta_{k+1})A^{(k)}+\theta_{k+1}B^{(k+1),\eta}$; the operators $G_k(\Omega_C,\cdot)$, $\Delta^{(k)}(\Omega_C,\cdot)$, $C^{(k)}_\Lambda(\Omega_C,\cdot)$ depend on the field only through the bonds of $\Omega_C$ \cite[(1.3)--(1.4)]{Ba4}. So the replacement of \cite[Prop.~2.6]{Ba2} is the pure domain change $\Omega_A\to\Omega_C$: its interaction-term half is \cref{thm:a07:dom}, and its quadratic-form half is the following. Layer distances \cite[(2.8)]{Ba2}: $\dist(\Lambda^{(k)}_{i+j},\Lambda^{(k)c}_i)\ge j(r-1)$ and $\dist(\Lambda^{(k-1)'}_7,\Lambda^{(k-1)'c}_2)\ge r-1$ ($L=5$). Standing lower bound: $R\ge R_0+3$, so that $r-2\ge R_0+1$ ($R$ may be enlarged, \cite[(2.7)]{Ba2}).

\begin{lemma}[discrete trace inequality]\label{lem:a06f:T}
Let $u$ be a function on a box $Y$ of $\eta\Z^d$ with all its bonds, $z$ a site, $e$ a lattice direction and $S_z\defeq\{z+m\eta e:0\le m<\eta^{-1}\}\subset Y$. Then $|u(z)|^2\le2\sum_{v\in S_z}\eta|u(v)|^2+2\sum_{b'\subset S_z}\eta|(D_Au)(b')|^2$.
\end{lemma}

\begin{proof}
$|U|=1$ gives $||u(b'_+)|-|u(b'_-)||\le\eta|D_Au(b')|$, so $|u(z)|\le|u(v)|+\sum_{b'\subset[z,v]}\eta|D_Au(b')|$ for every $v\in S_z$; average over $v$ with weight $\eta$ (total weight $1$), apply Cauchy--Schwarz and square.
\end{proof}

\begin{lemma}[domain difference of $\Delta^{(k)}$, decaying from the new boundary]\label{lem:a06f:Q2}
Let $\Omega'\subset\Omega$ be unions of big blocks, $A$ regular on $\Omega$, $\Sigma$ the set of bonds of $\Omega$ with exactly one end in $\Omega'$, and $\dist(\Sigma,\Omega^c)\ge R_0+1$. Then for $y,y'\in\Omega'^{(k)}$: $|\Delta^{(k)}(\Omega',A;y,y')-\Delta^{(k)}(\Omega,A;y,y')|\le c_Q\exp[-\delta_Q(|y-y'|+\dist(B(y),\Sigma)+\dist(B(y'),\Sigma))]$, $\delta_Q\defeq\min(\delta_0,\theta_0)/4$, $c_Q$ depending on $(c_0,\delta_0,\gamma_\emptyset,\theta_0,a,d)$. (This is stronger than \cite[(2.29)]{Ba1}, which decays from all of $\Omega'^c$.)
\end{lemma}

\begin{proof}
(1) With $H_\Omega\defeq-\Delta^{\eta,N}_{A,\Omega}+m^2t^2+a_kP_k(A)$: $H_\Omega=(H_{\Omega'}\oplus H_{\Omega\setminus\Omega'})+D_\Sigma^*D_\Sigma$, $D_\Sigma$ the covariant difference on the bonds of $\Sigma$ ($P_k$ and the mass are block diagonal and both pieces are unions of blocks). For $f'$ on $\Omega'$, $u'\defeq G_k(\Omega')f'$ and $u^-\defeq u'\oplus0$: $H_\Omega u^-=f'+D_\Sigma^*D_\Sigma u^-$, hence $(G_k(\Omega)f')|_{\Omega'}-u'=-(G_k(\Omega)D_\Sigma^*D_\Sigma u^-)|_{\Omega'}$. (2) With $q(y)\defeq Q_k^*(A)\delta_y$: $\delta\Delta(y,y')=a_k^2\ip{q(y)}{(G_k(\Omega)-G_k(\Omega'))q(y')}=-a_k^2\ip{D_\Sigma G_k(\Omega)q(y)}{D_\Sigma u^-}$, $u'=G_k(\Omega')q(y')$. (3) For $b=(z,w)\in\Sigma$ with $z\in\Omega'$, $|(D_\Sigma u^-)(b)|=\eta^{-1}|u'(z)|$. (4) By \cite[(2.25)]{Ba1} at $b$, the restricted form of the theorem of \cite{Ba4}, on the torus by \cref{lem:a06A:torus} (its hypotheses hold: $\Omega$ is a union of big blocks, $A$ is regular on $\Omega$, $\norm{q(y)}_\infty=1$, and $\dist(b,\Omega^c)\ge R_0$), $|D_bG_k(\Omega)q(y)|\le c_0e^{-\delta_0\dist(b,B(y))}$. (5) For a unit block $Y$ of $\Omega'$ let $\Sigma_Y$ be the bonds of $\Sigma$ with $z_b\in Y$ ($\#\Sigma_Y\le2d\eta^{-(d-1)}$); each $z_b$ lies on a face of $Y$, the segments away from $w$ stay in $Y$ and those from one face partition $Y$, so \cref{lem:a06f:T} over the $\le2d$ faces gives $\sum_{b\in\Sigma_Y}\eta^{d-1}|u'(z_b)|^2\le4d(\norm{\mathbf 1_Yu'}^2+\norm{\mathbf 1_YDu'}^2)$. (6) \Cref{lem:a06A:CT} on $\Omega'$ with $T=H_{\Omega'}\ge\gamma_\emptyset$ (\cref{lem:a06A:coer} with $\Lambda=\emptyset$; $\Omega'$ is a union of big blocks and $A$ is regular on it) gives $\norm{\mathbf 1_Yu'}\le2\gamma_\emptyset^{-1}e^{-\theta_0\dist(Y,B(y'))}$ and $\norm{\mathbf 1_YDu'}\le Ce^{-\theta_0\dist(Y,B(y'))}$, no depth restriction. (7) Cauchy--Schwarz on $\Sigma_Y$ with weights $\eta^{d-1}$: $|\delta\Delta|\le a^2\sum_Y\sum_{\Sigma_Y}\eta^d|D_bGq(y)|\eta^{-1}|u'(z_b)|\le C\sum_Ye^{-\delta_0\dist(Y,B(y))-\theta_0\dist(Y,B(y'))}$ over the blocks $Y$ adjacent to $\Sigma$. (8) For such $Y$, $\dist(Y,B(y))+\dist(Y,B(y'))\ge\frac12(\dist(B(y),B(y'))+\dist(B(y),\Sigma)+\dist(B(y'),\Sigma))-2$; keep half of the exponent for the sum over $Y$ and use $|y-y'|\le\dist(B(y),B(y'))+1$.
\end{proof}

\begin{lemma}[{\cite[(2.108)]{Ba2} is an identity}]\label{lem:a06f:C53}
For $\varphi$ on $\Lambda^{(k-1)'}_6$ put $\varphi_o\defeq(\Lambda^{(k-1)'}_6\cap\Lambda^{(k)c}_3)\varphi$, $\varphi_i\defeq\Lambda^{(k)}_3\varphi$, $\varphi_e\defeq(\Lambda^{(k)}_3\cap\Lambda^{(k)c}_4)\varphi$, $\varphi_4\defeq\Lambda^{(k)}_4\varphi$ ($\varphi_i=\varphi_e+\varphi_4$), $\Delta_A\defeq\Delta^{(k)}(\Omega_A,\tilde B)$, $\Delta_B\defeq\Delta^{(k)}(\Omega_B,\tilde B)$, $\Delta_C\defeq\Delta^{(k)}(\Omega_C,B^{(k+1),\eta})=\Delta^{(k)}(\Omega_C,\tilde B)$. Then $\frac12\ip{\Lambda'_6\varphi}{\Delta_A\Lambda'_6\varphi}=\frac12\ip{\varphi_o}{\Delta_B\varphi_o}+\ip{\varphi_o}{\Delta_B\varphi_e}+\frac12\ip{\varphi_i}{\Delta_C\varphi_i}+\frac12\ip{\Lambda'_6\varphi}{H_k\Lambda'_6\varphi}$ with $\frac12\ip\varphi{H_k\varphi}\defeq\frac12\ip{\varphi_o}{(\Delta_A-\Delta_B)\varphi_o}+\ip{\varphi_o}{(\Delta_A-\Delta_B)\varphi_e}+\ip{\varphi_o}{\Delta_A\varphi_4}+\frac12\ip{\varphi_i}{(\Delta_A-\Delta_C)\varphi_i}$. With the block-spin term $\frac12aL^{d-2}\sum|\psi-Q(\theta_k\tilde B)\varphi|^2$, unchanged on both sides, this is \cite[(2.108)]{Ba2} with the factor $\frac12$ of its fourth term restored.
\end{lemma}

\begin{proof}
$\Lambda^{(k)}_3\subset\Lambda^{(k)}_0\subset\Lambda^{(k-1)'}_7\subset\Lambda^{(k-1)'}_6$, so $\Lambda'_6\varphi=\varphi_o+\varphi_i$; the $\Delta$'s are real symmetric ($U\in SO(N)$, $q$ real antisymmetric); expand and subtract. All forms are defined: $B^k(\Lambda'_6\cap\Lambda^c_3)$ and $B^k(\Lambda_3\cap\Lambda_4^c)$ lie in $\Omega_B$, $B^k(\Lambda_3)$ in $\Omega_C$.
\end{proof}

\begin{lemma}[{\cite[(2.109)]{Ba2}}]\label{lem:a06f:C54}
For $y,y'\in\Lambda^{(k-1)'}_6$ the kernel of $H_k$ obeys $|h_k(y,y')|\le C_He^{-\delta_H(r-2)}e^{-\delta_H|y-y'|}$, $\delta_H\defeq\min(\delta_Q,\delta_7/2)/2$ ($\delta_7$ the rate of \cref{prop:a06A:Ba1P22}, the form of \cite[(2.27)]{Ba1} proved in Step~3(A)); and under \cite[(2.55)]{Ba2}, for $t\le t_0(\kappa)$, $\frac12|\ip{\Lambda'_6\varphi}{H_k\Lambda'_6\varphi}|\le Cc_1^2p(t/L)^2\lambda(t/L)^{-1/2}e^{-(\delta_H/2)(r-2)}|\Lambda^{(k)}_3|\le C_\kappa t^\kappa|\Lambda^{(k)}_3|$, the printed error of \cite[(2.108)]{Ba2}.
\end{lemma}

\begin{proof}
Four pieces. (oo), (oe): \cref{lem:a06f:Q2} with $(\Omega,\Omega')=(\Omega_A,\Omega_B)$; $\Omega_A\setminus\Omega_B=B^k(\Lambda^{(k)}_5)$, so $\Sigma$ consists of the bonds crossing $\partial B^k(\Lambda^{(k)}_5)$, with $\dist(\Sigma,\Omega_A^c)\ge\dist(\Lambda^{(k)}_5,\Lambda^{(k-1)'c}_2)-1\ge r-2\ge R_0+1$; for $y\in\Lambda_3^c$, $\dist(B(y),\Sigma)\ge2(r-1)-2$, for $y'\in\Lambda_3\cap\Lambda_4^c$, $\dist(B(y'),\Sigma)\ge r-2$; so $|(\Delta_A-\Delta_B)(y,y')|\le c_Qe^{-(\delta_Q/2)(r-2)}e^{-(\delta_Q/2)(|y-y'|+\dist(B(y),\Sigma)+\dist(B(y'),\Sigma))}$. (o4): for $y\in\Lambda_3^c$, $y'\in\Lambda_4$, $|y-y'|\ge r-1$, and \cref{prop:a06A:Ba1P22} (\cite[(2.27)]{Ba1}, under the hypotheses of \cref{lem:a06f:Q2} on $\Omega_A$) gives $|\Delta_A(y,y')|\le c_7e^{-(\delta_7/2)(r-1)}e^{-(\delta_7/2)|y-y'|}$. (ii): \cref{lem:a06f:Q2} with $(\Omega_A,\Omega_C)$, $\Sigma$ the bonds crossing $\partial\Omega_C$, $\dist(\Sigma,\Omega_A^c)\ge r-2$, and $\dist(B(y),\Sigma)\ge r-2$ for $y\in\Lambda_3$. Assembling the kernels gives the kernel bound. Error: $\sup|\varphi|\le c_1p(t/L)\lambda(t/L)^{-1/4}$ on $\Lambda'_6\subset\Lambda'_{-1}$ by \cite[(2.55)]{Ba2}; the (oo) sum over $y,y'\in\Lambda_3^c$ is $\le Ce^{-(\delta_Q/2)(r-2)}\sum_ye^{-(\delta_Q/4)\dist(B(y),\Sigma)}\le Ce^{-(\delta_Q/2)(r-2)}|\Lambda^{(k)}_5|$ (at most $|\Lambda_5|$ unit blocks are adjacent to $\Sigma$ from inside); in (oe), (o4), (ii) the second index runs over $\Lambda_3$ or $\Lambda_4$, giving $\le C|\Lambda^{(k)}_3|$ times the exponential; then $|\Lambda_5|\le|\Lambda_3|$, $p(t/L)\le(1+\ln L)^{\bar p}p(t)$, $\lambda(t/L)=\lambda_k/L$, and $p^2\lambda^{-1/2}e^{-\delta r}\le t^\kappa$ for $t\le t_0(\kappa)$.
\end{proof}

With \cite[(2.29)]{Ba1} in place of \cref{lem:a06f:Q2} the error would be $O(t^\kappa)|\Lambda^{(k-1)'}_6|$, still summable at $\kappa>d$ but not the printed $|\Lambda^{(k)}_3|$.

\subsubsection*{From \cite[(3.5)]{Ba2} at $k=0$ to \cite[(3.21)]{Ba2}, and the payment of \cite[Sect.~3C]{Ba2}}

\begin{lemma}[monotone nesting]\label{lem:a06f:C141}
Every renormalisation transformation $T^t_{a,L}$, $T^t_{a,L,\tilde A}$ of \cite[Sect.~3]{Ba2} is $F\mapsto\int t(\cdot,F\text{-arguments})F$ with a kernel $t\ge0$ normalised by \cite[(3.32)]{Ba2}; hence $F\le F'$ pointwise implies $T[F]\le T[F']$, and nested brackets preserve inequalities.
\end{lemma}

\begin{lemma}[{\cite[(3.21)]{Ba2}}]\label{lem:a06f:C142}
\cite[(3.5)]{Ba2} holds at $k=0$ (\cref{prop:a06d:C211}, the finite downward induction from the base case). At $k=0$ the innermost bracket of \cite[(3.5)]{Ba2} is the bare Gaussian density $\exp(-\frac12\ip{A_0}{(-\Delta^\eps+\lmu^2)A_0})\exp(-\frac12\ip{\phi_0}{(-\Delta^\eps_{\tilde A^\eps}+m^2)\phi_0})$ with the bare characteristic functions \cite[(2.34)--(2.42), (2.47)--(2.49)]{Ba2}; each characteristic function factorises into a vector condition and a scalar condition with the vector configuration as parameter, $\zeta$ placed inside the scalar bracket; integrating first over the scalar fields at fixed vector fields (Tonelli) and \cref{lem:a06f:C141} give \cite[(3.21)]{Ba2}, whose last factor is $\exp(\sum_{k=1}^KC_\kappa t_k^{\kappa_0}|T_\eps|)\le\exp(C_\kappa\eps_0^{\kappa_0}(1-L^{-\kappa_0})^{-1}|T_\eps|)$, i.e.\ \cite[(2.118)]{Ba2}.
\end{lemma}

\begin{lemma}[{the residue of \cite[Sect.~3B]{Ba2}}]\label{lem:a06f:C143}
In \cite[(3.30)--(3.40)]{Ba2} the only factors left besides $\prod\zeta''$ and $\exp(E_0)$ are: from \cite[(3.31)]{Ba2}, $\exp(c_{3.29}t_k^{\kappa_0}|\Lambda_k|)$ with $c_{3.29}$ the remainder constant of \cite[Prop.~3.1, (3.29)]{Ba2} (\cref{thm:a06:BALA}, the $\eps_0$-freedom paragraph of Step~3(C$_2$)); from \cite[(3.33)]{Ba2}, $2^{N/2}\exp(-\frac14ap(t_k)^2)\le\exp(-\frac18ap(t_k)^2)$ per exceptional block, valid once $b_0^2\ge4N\ln2/a$ (a condition on $b_0$, fixed before $\eps_0$); from \cite[(3.36), (3.40)]{Ba2}, the Jacobians $\exp(\frac N2\ln2\sum_k|\Lambda_k|)$ and $\exp(\frac d2\ln2\sum_k|\Lambda_k|)$; and from \cite[(3.38)]{Ba2} the diamagnetic inequality and the normalisations \cite[(3.32)]{Ba2} (equalities). Here $\Lambda_k\defeq\Lambda^{(k-1)'}_5\cap\Lambda^{(k)c}_5$, counted in points of the scale-$k$ unit lattice. So the $O(1)$ of \cite[(3.39)--(3.40)]{Ba2} is $c_5\defeq\frac{N+d}2\ln2+c_{3.29}\eps_0^{\kappa_0}\le\frac{N+d}2\ln2+c_{3.29}$, and the $O(1)$ in front of $(L^k\eps)^{\kappa_0}|\Lambda^{(k-1)'}_7\cap\Lambda^{(k)c}_7|$ in \cite[(2.117), (3.42)]{Ba2} is $c_7\defeq C_{2.57}+C_{2.78}$ (Step~3(C$_2$), the $O(1)$ lines). \cite[(3.40)]{Ba2} prints $\sum_{k=1}^K|\Lambda_k|$ where \cite[(3.36), (3.42)]{Ba2} print $\sum_{k=0}^K$; the constants use $\sum_{k=0}^K$.
\end{lemma}

\begin{lemma}[the payment]\label{lem:a06f:C144}
The residue $\exp(c_5\sum_k|\Lambda_k|)$ of \cref{lem:a06f:C143} is exactly the second factor of the left side of \cite[(3.42)]{Ba2} with $O(1)=c_5$, and the first factor carries $O(1)=c_7$. So \cite[(2.117)]{Ba2}, hence the upper bound, follows from \cite[(3.42)]{Ba2} with the constants $(c_5,c_7)$, which are free of $\eps$, $K$, the torus, the history and $\eps_0$; \cite[(3.42)]{Ba2} itself is \cref{thm:a06f:R7} below, which imposes its own cap on $\eps_0$ after $c_5$, $c_7$ are fixed.
\end{lemma}

\begin{proof}[Proof of \cref{lem:a06f:C141,lem:a06f:C142,lem:a06f:C143,lem:a06f:C144}]
\Cref{lem:a06f:C141}: positivity of the Gaussian kernels. \Cref{lem:a06f:C142}: as stated, the $k=0$ bracket being the bare density by the definitions of the first step (\cite[p.~588]{Ba2}: ``using inequality (3.5) for $k=0$''), and $t_k=L^{k-K}t_K\le L^{k-K}\eps_0$. \Cref{lem:a06f:C143}: arithmetic on the printed steps, with $p(t)\ge b_0$ for the $b_0$-condition. \Cref{lem:a06f:C144}: $|\Lambda_k|=|\Lambda^{(k-1)'}_5\cap\Lambda^{(k)c}_5|$ by definition (\cite[p.~592]{Ba2}: ``$|\cdot|$ means the number of elements''), an identity of terms.
\end{proof}
